	Dans cet exercice, on considère la figure ci-contre.

	\medskip

	\begin{minipage}{0.47\linewidth}
Les points A, B et M sont alignés.

Les points A, C et N sont alignés.

Le triangle ABC est rectangle en B.

Le triangle AMN est rectangle en M.
\end{minipage}\hfill
\begin{minipage}{0.49\linewidth}
%début PSTricks
\psset{unit=1cm,arrowsize=2pt 3}
\begin{pspicture}(7,4.8)
	%\psgrid
	\psline(0,0.6)(7,4.7)
	\psline(1,1.2)(6,1.2)(6,4.1)%AMN
	\psline(4,1.2)(4,2.96)%BC
	\uput[dl](1,1.2){A} \uput[r](6,1.2){M} \uput[dr](6,4.1){N}
	\uput[ur](4,1.2){B} \uput[ur](3.9,3.13){C}
	\psline[linewidth=0.8pt]{<->}(1,0.8)(4,0.8)\uput[d](2.5,0.8){6,4 cm}
	\psline[linewidth=0.8pt]{<->}(4,0.8)(6,0.8)\uput[d](5,0.8){3,2 cm}
	\psline[linewidth=0.8pt]{<->}(0.9,1.3)(3.95,3.11)\rput{31}(2.5,2.45){8 cm}
	\psframe(4,1.2)(3.75,1.45)
	\psline(6,1.2)(6,0)%M..
\end{pspicture}
%Fin PSTricks
%Tikz
%\begin{center}
%\begin{tikzpicture}[scale=0.55]
%\coordinate (A) at (0,0);
%\coordinate (B) at (6.4,0);
%\coordinate (M) at (9.6,0);
%\coordinate (C) at (6.4,4.8);
%\coordinate (N) at (9.6,7.2);
%\draw (A) -- (M);
%\draw (A) -- (N);
%\draw (B) -- (C);
%\draw (M) -- (N);
%\draw (5.9,0) rectangle (6.4,0.5);
%\draw (9.1,0) rectangle (9.6,0.5);
%\node[below] at (A) {$A$};
%\node[below] at (B) {$B$};
%\node[below right] at (M) {M};
%\node[above left] at (C) {C};
%\node[above right] at (N) {N};
%\node[below] at (3.2,0) {6,4 cm};
%\node[below] at (8,0) {3,2 cm};
%\node[left] at (3.2,2.4) {8 cm};
%\end{tikzpicture}
%\end{center}
\end{minipage}

\medskip

On donne :

AB $= 6,4$ cm ; BM $= 3,2$ cm et AC $= 8$ cm.

\emph{La figure n'est pas en vraie grandeur.}

\begin{enumerate}
\item Tracer sur la copie le triangle ABC en vraie grandeur et en laissant les traits de construction.
\item Démontrer que BC $= 4,8$ cm.
\item Justifier que les droites (BC) et (MN) sont parallèles.
\item Démontrer que MN $= 7,2$ cm et AN $= 12$ cm.
\item Le périmètre du triangle ABC et le périmètre du quadrilatère BMNC sont-ils égaux ?\\
\textbf{Argumenter la réponse en précisant la démarche.}
\end{enumerate}
